% Documento de trabajo: análisis para un posible capítulo.
% No contiene prosa destinada al manuscrito.

\section*{La bolsa de los vientos: líneas directrices}

\subsection*{Decisión estructural}

En la versión destinada a Jot Down, «Circe» puede mantenerse como una sola
entrega: el medio digital admite su extensión y la sucesión de episodios
funciona como narración continua.

En el libro conviene valorar una división entre «Cíclope» y «Circe». El nuevo
capítulo se titularía «La bolsa de los vientos» y reuniría el episodio de Eolo
y, probablemente, el de los lestrigones. «Circe» comenzaría directamente con
la llegada a Eea.

Los lotófagos quedan fuera. Su episodio es anterior al del cíclope y su
inclusión obligaría a romper la progresión cronológica. Tampoco es necesario
para sostener el argumento del nuevo capítulo.

\subsection*{Función dentro del libro}

La división permitiría separar dos movimientos distintos:

\begin{itemize}
\item Eolo y los lestrigones: fracaso del regreso, desconfianza y destrucción
      de la flota.
\item Circe: transformación, deseo, tejido, canto, culpa y redención.
\end{itemize}

Al retirar la larga antesala de Eolo y los lestrigones, Circe aparecería antes
y su capítulo ganaría concentración. El nuevo capítulo, más breve, introduciría
además una variación de ritmo entre «Cíclope» y «Circe».

\subsection*{Tesis central}

La bolsa no debe tratarse únicamente como un objeto maravilloso. Puede
funcionar como símbolo de la fractura interna de la expedición:

\begin{itemize}
\item Eolo entrega a Odiseo la posibilidad material de regresar.
\item Odiseo conoce el contenido de la bolsa, pero no confía la vigilancia a
      sus hombres ni parece explicarles qué contiene.
\item La tripulación sospecha que el capitán esconde oro y que vuelve a
      reservarse la mejor parte del botín.
\item Odiseo no confía en sus hombres y sus hombres no confían en Odiseo.
\item La catástrofe no procede directamente de Poseidón ni de una maldición,
      sino del secreto, la envidia, el agotamiento y la estupidez.
\end{itemize}

El episodio ofrece así una reflexión sobre el fracaso del liderazgo, no solo
sobre la indisciplina de la tripulación. La conducta de los hombres es
desastrosa, pero la decisión de Odiseo de monopolizar la información y el
control de la bolsa forma parte del problema.

\subsection*{Ítaca a la vista}

La fuerza dramática del episodio reside en que la patria llega a verse. Las
hogueras de Ítaca dejan de ser un recuerdo o una esperanza: están físicamente
al alcance de los navegantes. La bolsa se abre precisamente cuando el regreso
parece consumado.

Debe darse relevancia al despertar de Odiseo. Al ver que Ítaca vuelve a
alejarse, considera arrojarse al mar y tiene que escoger entre morir o soportar
que la vida continúe. Es uno de los momentos en que el héroe aparece más
desnudo y menos protegido por sus tretas.

\subsection*{Eolo y la interpretación religiosa}

Al regresar a Eolia, Eolo interpreta la desgracia como prueba de que Odiseo es
aborrecido por los dioses. Esa explicación es razonable dentro de su mundo,
pero el lector acaba de presenciar causas enteramente humanas.

El contraste permite señalar cómo los personajes atribuyen a los dioses las
consecuencias de sus propios errores. También introduce un límite de la
\emph{xenía}: Eolo recibe y agasaja al extranjero afortunado, pero expulsa al
que regresa derrotado porque teme ayudar a un hombre aparentemente maldito.

Eolo no debe ser llamado dios. Homero lo presenta como hijo de Hipotas,
querido por los inmortales y designado por Zeus custodio o señor de los
vientos.

\subsection*{Función de los lestrigones}

Los lestrigones pueden adquirir mayor importancia que en la versión actual.
No son únicamente una repetición colectiva de Polifemo:

\begin{itemize}
\item Polifemo mata a varios compañeros; los lestrigones destruyen la flota.
\item El puerto, aparentemente seguro, se convierte en una trampa geográfica.
\item Once naves entran en la ensenada y quedan encerradas.
\item La nave de Odiseo permanece fuera y es la única que puede escapar.
\item La separación puede interpretarse como prudencia, desconfianza o hábito
      del caudillo de mantenerse aparte de los hombres que dirige.
\item Los marineros son ensartados como peces y las naves quedan destrozadas
      por las rocas arrojadas desde los acantilados.
\end{itemize}

La apertura de la bolsa destruye la posibilidad inmediata del regreso. Los
lestrigones destruyen después la expedición misma. Al escapar, Odiseo ya no es
el jefe de una flota, sino el capitán de una única nave llena de supervivientes.
Este cambio de escala prepara psicológica y narrativamente la llegada a Circe.

\subsection*{Material que habría que trasladar desde «Circe»}

Sin eliminarlo hasta que la división sea definitiva:

\begin{itemize}
\item presentación de Eolia y de la familia de Eolo;
\item hospitalidad durante un mes;
\item entrega de la bolsa;
\item navegación durante nueve días e Ítaca a la vista;
\item conversación de los compañeros y apertura de la bolsa;
\item desesperación de Odiseo;
\item regreso a Eolia y expulsión;
\item navegación a remo;
\item presentación de los lestrigones;
\item destrucción de las once naves y huida de la última.
\end{itemize}

Para ampliar a los lestrigones habría que recuperar del canto X:

\begin{itemize}
\item la descripción del puerto cerrado por farallones;
\item la decisión de Odiseo de dejar su nave fuera;
\item la expedición de reconocimiento;
\item el encuentro con la hija de Antífates;
\item el primer hombre devorado;
\item el ataque desde las alturas;
\item la comparación de los marineros con peces ensartados.
\end{itemize}

\subsection*{Enlace con «Circe»}

El nuevo capítulo debería terminar con la única nave superviviente rumbo a
Eea. «Circe» podría comenzar con los versos que presentan la isla y a la diosa
de las bellas trenzas. De este modo, la protagonista aparecería al comienzo de
su propio capítulo y la pérdida de la flota explicaría el estado de agotamiento,
temor y desconfianza con el que los aqueos llegan a su morada.
